\section{Objetivos}
\subsection{Objetivo general}
{\justify Desarrollar e implementar un plugin de localización en tiempo real mediante sensores Estimote con tecnología Ultra-Wideband (UWB), integrado a Unity, como parte de una aplicación de realidad aumentada que facilite los recorridos virtuales dentro del Centro de Innovación y Tecnología (CIT) del campus de la Universidad del Valle de Guatemala.}

\subsection{Objetivos específicos}
\begin{enumerate}[label=\thesubsection.\arabic*]
    \item Diseñar e implementar un algoritmo de trilateración
para estimar la posición bidimensional del usuario dentro del campus universitario,
mediante el procesamiento de las distancias medidas hacia múltiples sensores Estimote UWB distribuidos en el entorno.
    \item Integrar el SDK nativo de Estimote para iOS con Unity
para permitir la obtención y transferencia en tiempo real de datos de localización hacia la aplicación de realidad aumentada,
mediante el desarrollo de un plugin nativo en Swift que sirva de puente entre ambas plataformas.
    \item Aplicar técnicas de filtrado de señales
para mejorar la estabilidad del sistema de localización y reducir el efecto de ruido o jitter en las mediciones,
mediante la combinación de los datos obtenidos por trilateración y los registros del acelerómetro del dispositivo.
    \item Documentar el desarrollo del sistema de localización
para asegurar su mantenibilidad, facilitar su futura extensión a dispositivos Android y permitir su continuidad,
mediante la elaboración de documentación técnica clara sobre la arquitectura del plugin, las decisiones de diseño y los procesos de integración.
    \item Evaluar cualitativamente la precisión y suavidad del sistema de localización
para validar su funcionamiento en condiciones reales y su aporte a la experiencia de usuario dentro de la aplicación,
mediante la observación directa del comportamiento del sistema y el análisis de las distribuciones de datos de posición obtenidas durante pruebas manuales.       
\end{enumerate}

\newpage